\subsection{Proved conclusions and pending results}
The settled results are the two bounded simulations (Theorems~\ref{eqv:forward} and \ref{eqv:reverse}), the conditional approximation theorem and its effective rotation construction (Theorems~\ref{eqv:approx} and \ref{eqv:rotations}), structured unitary soundness (Theorem~\ref{thm:sA-sound}), exact and recursive restoration tests (Theorems~\ref{thm:sA-clean} and \ref{thm:sA-recursive}), and the sufficient product-output restoration condition with its explicit limitations (Theorems~\ref{thm:sA-classical} and \ref{thm:sA-four}). Universal auxiliary restoration is false, as witnessed by the four-qubit circuit.
Counterexample~\ref{ex:sA-rec-limits} limits the reset tests themselves:
containing subspaces lose necessity, final cleanliness does not imply
cleanliness of inner marks, marginal restoration does not license
deletion, and a unitary reduced channel need not reset the auxiliary
register. Proposition~\ref{eqv:restrictions} and
Counterexamples~\ref{eqv:exact-obstruction} and
\ref{eqv:phase-obstruction} prove that the capacity, stopping, exact-angle
and phase restrictions cannot be suppressed;
Corollary~\ref{eqv:adaptive-circuits} packages the two simulations for
bounded adaptive device data. The resource and conditional bounded-interpreter results below are proved; compiling general dynamic source mark/rewind into a certified interpreter remains open. Open questions and exclusions are collected in Section~\ref{sec:open-problems}.

Primitive~10 has the controlled-reversible meaning of
Lemma~\ref{lem:read-product} when the subsequent operations preserve the
read register's computational-basis sectors. A gate changing that register
removes this guarantee; the full joint unitary and its retained records
then specify the semantics.

Theorem~\ref{thm:resA-resources} gives exact gate and serial-time counts for the prescribed expansion, sufficient qubit counts, and an explicit compiler bound. A mark doubles its body's work and adds its keep cost. Repeating this nesting gives a tight exponential count for that expansion. This is not a lower bound for optimal equivalent circuits: with an empty keep the expanded circuit can already denote identity. Repeated inverse uses reuse the same named records; extra gates alone do not require extra logical wires.

Theorem~\ref{thm:resA-delayed} proves a restricted alternative: retain all evaluator outputs and private wires until the final copy, then invert the evaluators in reverse order. For uniformly positive bounded evaluator costs and bounded copy costs, its serial time is linear in the number of levels, while the prescribed eager construction is exponential. The stated sufficient space allocations need not be strictly separated. More generally, keeping an inner record until the outer rewind and deleting its immediate inverse preserves the whole retained operator exactly when the inner forward operation commutes with the outer keep conjugated by the later computation (Theorem~\ref{thm:resA-deletion}). A weaker exact test applies on a specified input embedding. A clean inner computation alone does not suffice: the five-qubit Counterexample~\ref{ex:resA-quantum} changes the probability of $r=1$ from $1/2$ to zero, even after the machine register is traced. Certified repeated deletions give linear time in the once-counted body size only under Corollary~\ref{cor:resA-flat}'s tests and restrictions.

For Tier~B, Theorem~\ref{thm:tB-sound} proves that a bounded interpreter run is unitary on all retained logical registers when the supplied certificate verifies the stated rules, completions and reached subspaces. Under the whole-sector matching encodings, each undo with its label routing is the exact adjoint of its completed forward operation, including its retained records. A supplied clean implementation of the whole tick returns workspace to zero at every tick. Theorem~\ref{thm:tB-A} realizes Tier~A as a static specialization of the clock schedule with all its machine registers retained; it constructs no dynamic certificate and implies no machine reset.

Theorem~\ref{thm:tB-clean} gives an exact synchronization and disposal test. For all data inputs and a fixed pure initial machine state, tracing the interpreter register leaves a unitary data channel exactly when the entire final interpreter register factors in one fixed pure state independent of the input. This includes position, clock, journal, stack and routing records. The fixed final vector need not be the initial one: disposal and logical reset remain distinct. Equal durations alone do not establish the test. In the four-qubit Counterexample~\ref{ex:tB-delay}, exact undo restores the target after two or three ticks according to the data value. At the fixed three-tick endpoint the local timing record remains correlated with the data; disposing of it sends the control qubit's $\ket+$ state to the equal mixture of $\ket0$ and $\ket1$. This is a finite matrix example, not an exhibited dynamic-certificate instance.

Under the supplied-circuit and constant-access hypotheses of Definition~\ref{def:tB-cost}, Theorem~\ref{thm:tB-cost} gives the per-tick gate cost
\[
 O\bigl(\ell(a+w)+N(u+w)+d(v+w)+\beta+g+b+w+1\bigr).
\]
Here $\ell,N,d$ are instruction-table, journal and mark-stack capacities; $a,u,v$ are their address widths, $w$ the cell width, $\beta$ the supplied Boolean-circuit cost including inverse routing, $g$ the supplied action-circuit cost, and $b$ the clock width. The first three terms charge linear scans; there is no assumed memory-access oracle. Exactly $T$ copies of the tick circuit give the total gate cost; the sufficient wire allocation includes all retained records and clean workspace. Proposition~\ref{prop:tB-nesting} leaves nesting dependence in these parameters and the horizon: a journal alone does not remove the prescribed expansion's doubling, and capacity $N=O(T)$ gives a quadratic-in-$T$ scan term with its width factors retained.

The open step is compiling general dynamic source mark/rewind into an interpreter meeting the certificate, particularly data-dependent repeated mark visits and an inverse stopped at a coherent mark identifier (Open problem~\ref{prob:tB-compiler}). No such general compiler, source-level synchronization condition under arbitrary gates, unbounded-runtime extension or executor-equivalence theorem is proved.

\subsection{Every restriction of the admitted program class}
\begin{enumerate}
\item All instances and register lists are finite and fixed independently of the unknown quantum input. Uniform families are generated by one classical Turing machine. Polynomial uniformity additionally requires polynomial numerical bounds and polynomial generation time. Compilation counts charge the supplied descriptions, binary addresses and explicit output lists; numerical bounds in those estimates are supplied in unary. A binary horizon still costs its numerical number of execution steps.
\item Each device supplies the complete field layout, validity test, initial configuration, canonical total decoder, total update and total output circuits, including invalid-code outputs and fault behavior. All live-wire checks and flag-preservation rules hold on every configuration code. Unspecified arithmetic, addressing, overflow, division and call conventions are not identified with any particular source language. No efficiency of an arbitrary supplied Boolean function is presumed.
\item The initial configuration is valid and its live flags identify exactly the input wires. Allocations select never-used zero factors. Names are never reused, dropped names remain inaccessible, and no later operation acts on them. There is no device storage-release instruction. The output quantum list is fixed and live at final time for every reachable outcome list; the entire complementary retained register is traced.
\item Observed protocols for primitives 1--9 contain no marks and exclude coherent-read mode on every code. They use an input-independent horizon of at least one tick, terminal padding and a separate final timeout snapshot within that horizon. Each forward tick has distinct blank sample and configuration records. Measurements are computational-basis observations; classical output is a specified final tape or register. Equivalence concerns the final instrument on the same input and output lists, with arbitrary references; it preserves neither write times nor stopping times and describes no intermediate interaction.
\item A reversible Turing-machine transition is partial and injective under the stated tape convention. Finite total unitary use fixes its completion. Agreement with the original machine requires each forward transition to stay within the valid finite window and restricted transition domain. Efficient device simulation uses supplied circuits and claims no efficient arbitrary completion.
\item A Tier~A fixed-space marked term uses total unitary atoms, total Boolean keeps and a fixed retained register list. Marks and their pairings are static. Calls or loops containing marks need a supplied fixed finite syntactic expansion; general recursion and data-dependent numbers of mark visits are excluded. Bounded loops without marks may instead use scheduled ticks. Coherent tick decoders exclude measurement and allocation on every code. Continuation addresses must fit their fields.
\item No allocation, drop, observed measurement, external console observation or actual stop occurs inside a marked body. An unobserved console tape and bounded idle statuses may be updated. Logical records and implementation workspace are allocated before each enclosing run. Source allocation inside a mark needs a separate zero-input reservation proof and is not admitted automatically.
\item Atom occurrences have distinct retained record names; reversal reuses those same names in reverse order. Any factor needed by an enclosing inverse is retained. Nonblank records still have a total unitary action, but ordinary forward instruction semantics requires blankness. Continuation is a separate recorded update; position restoration and machine reset are additional properties.
\item Each keep uses a supplied total Boolean circuit on the whole specified machine list and a fixed list of distinct target wires disjoint from that machine list. Targets lie in the enclosing context; there is no computed-address exception. Keep entries may overlap targets. A target used by the body is allowed, with no automatic restoration guarantee.
\item Exact finite-alphabet implementation claims require supplied exact clean circuits for the atoms. Implementation workspace starts and ends at zero; no equality on its nonzero states is claimed. All physical qubits, including logical records, are retained before output partial trace. An ideal-angle atom has no automatic exact finite-alphabet implementation.
\item A controlled menu has a finite fixed collection of base structured terms, a common payload list, and a disjoint selector with sufficient capacity. Payloads contain no further menus, preserve the selector, and supply actual-phase exact finite-alphabet circuit lists for the implementation claim. Records are distinct across entries and occurrences and are retained, with the selector, for enclosing inverses. This does not define dynamic source mark stacks.
\item Restoration tests apply on the specified input subspace. Exact prefix images give necessity and sufficiency for actual inputs; containing subspaces give sufficiency alone. Blank-record reuse also requires an incoming zero condition. Logical reset points treat menu atoms as indivisible; boundaries within their entries or within physical implementations require the specified full circuit, its own initial subspace, all workspace, and the selector where present. Marginal restoration or a successful reset does not permit deletion of a register still required by an inverse. Final cleanliness does not imply cleanliness of every inner mark.
\item The sufficient product-output restoration theorem requires a total unitary forward operation independent of the keep target. For each basis input its forward data and machine outputs must be a product with a basis-valued machine output, with all keep controls included in that machine list. It need not preserve the input basis value. Removing the basis-output or unread-target requirement invalidates the universal claim; no necessity of these sufficient hypotheses for individual clean operations is asserted. The four-qubit witness uses explicit copying atoms rather than full-snapshot source-read instructions.
\item Approximation uses rational requests, counts every expanded factor occurrence including inverses and controls, bounds actual matrices with fixed uniform length constants, and reverses the same approximating circuit for inverse uses. General certificate-generation time is charged separately and need not be polynomial. The discharged case is determinant-one one-qubit matrices with supplied dyadic entry approximations of stated precision and cost, with a fixed number of controls. Arbitrary scalar phases, other additional matrices and growing control counts are outside that synthesis theorem. No matrix-entry procedure is derived merely from an angle description.
\item All resource gate counts use the prescribed unsimplified lists over the adjoint-closed alphabet $\mathcal G^{\pm}$ and serial execution. Sufficient wire allocations are not minimum-space lower bounds. The exponential lower counts apply only to this expansion; there is no lower bound here for optimal equivalent circuits, parallel depth or optimal time with a space restriction. These resource results do not imply logical-record reset. Zero-initialisation
of non-input wires, if charged at one tick per wire, adds
$N(P)-q_{\rm in}$ ticks, and preparing a nonzero initial configuration
is additional (Theorem~\ref{thm:resA-resources}).
\item The snapshot wire and scratch bounds use one shared configuration width, sample width and pair of device-circuit bounds across all ticks. They remain parametric in declared extra logical wires, arbitrary supplied-atom scratch and keep-circuit sizes. No efficiency bound for an arbitrary total completion is supplied. Payload record occurrences and explicit selector computations are charged once in the named terms; an externally supplied selector contributes its retained width, and its earlier computation cost is outside the displayed allocation. The compiler bound assumes an uncompressed prefix syntax encoding and already supplied delimited gate lists; constructing those lists is excluded.
\item The delayed-evaluator comparison requires total local evaluators, zero output and private inputs, a preserved basis input, an exact basis output and a product private factor for each basis input. That private factor may be a quantum state. All evaluator implementation wires are included in the declared private widths. XOR-copy circuits use no scratch and act only on the output and a disjoint target. Final targets never become evaluator inputs; earlier private wires remain untouched until their inverses, and no intervening operation changes retained values. Equality is of initialized retained isometries, up to auxiliary padding and factor permutation, and need not hold on nonzero auxiliary inputs. The linear versus exponential time comparison also requires the stated uniform cost bounds.
\item Deleting an inner inverse requires the exact commutator condition for equality on every retained input, or the stated weaker condition for the fixed input embedding. It does not itself reset a record or justify reuse of logical zero cells. The repeated-deletion corollary excludes menu atoms, requires a valid test at every deletion and requires that these deletions reduce the body to one visit per specified occurrence. It retains every logical wire and keeps the supplied exact clean local circuits. It provides no general transformation for arbitrary intervening quantum gates and no intermediate or optimal time--space trade-off. The corollary as stated has a single outermost mark.
\item Tier~B fixes finite capacities, field widths, program and an input-independent horizon $T$ before the unknown input; all wires are reserved initially. Its interpretation requires a supplied certificate for a linear initialized subspace and every reached sector component, including post-keep and partial-undo states. Rules have separately verified orthogonal input and output sectors. Capacity failure or an invalid mark identifier on the reached subspace defeats certification of successful source execution. No observation, allocation, drop or physical stop occurs during the run.
\item Every dynamic completed step has explicit forward and undo payload encodings. The matching identities (B2a), whole-sector routing realizations and consecutive-entry identifications hold on the entire retained payload, including post-keep vectors. All control changes, visit identifiers, popped records and earlier rewinds are retained or reversibly cleared with their exact action. Repeated visits require distinct live identifiers or a proved reversible reuse rule; merely naming mark and rewind proves none of these obligations. Constructing this interface for general dynamic source code remains open.
\item Tier~B circuit claims supply the specified completion and a clean implementation of the whole tick; abstract finite-dimensional extension gives neither an efficient circuit nor exact finite-alphabet synthesis of an arbitrary prescribed completion. The per-tick bound assumes the stated microcode, an absolute constant number of journal/stack accesses, Boolean-circuit sizes including ports and inverse routing, and supplied action, other-register and scratch costs. The run uses exactly $T$ tick circuits even when some sector components are idle. No polynomial source bound for $T$ follows. The static Tier~A embedding is a logical-factor schedule and supplies no dynamic certificate or unit-cost interpretation of a compound factor.
\item The Tier~B disposal equivalence tests every data input from a fixed pure initial interpreter state and requires collective pure factorization. A source interpretation additionally needs the specified zero clock and workspace factors and certified initial-subspace containment. Equal completion times or constant individual mixed marginals do not suffice. The unequal-duration example demonstrates failure for its displayed total unitary; it asserts neither that every unequal-duration run decoheres nor that the example supplies a general source translation.
\end{enumerate}

\subsection{Every correction needed for primitives 1--9}
\begin{enumerate}
\item (primitives 1--6) Integer arithmetic, comparison, memory access, conditional jumps, calls, returns, stopping and console output need complete classical functions and finite capacities. Growing uniform families replace one fixed finite device when comparison with unbounded computation is intended. Word arithmetic is the supplied bounded-width function; larger integers require charged multiword storage and work. Console output is a tape with specified overflow and final extraction.
\item (all primitives) The instruction names do not determine behaviour on
illegal codes. Each device must supply total decoder, update and output
circuits in the canonical layout, including a specified fault configuration
and the output on every invalid encoding (Definition~\ref{def:device}).
\item (primitive 5; all primitives) A horizon, padded terminal states and a scheduled timeout are extra data. The nine instruction names alone imply no termination theorem. Truncating an unbounded execution and adding timeout can change its instrument.
\item (primitive 9) Later quantum choices may depend on observed outcomes. Retained measurement copies, total classical updates and final environment trace define this adaptivity precisely. The bounded equivalence corollary (Corollary~\ref{eqv:adaptive-circuits}) takes an already supplied device presentation; it does not silently compile an undefined controller notation.
\item (primitive 7) Fresh allocation means a zero factor that has never been used. Drop is simulated by retaining an inaccessible factor and tracing it at output; storage reclamation and name reuse are not supplied by that simulation.
\item (primitive 8) Exact gate claims use the actual $H,T,\operatorname{CNOT}$ alphabet. $T^\dagger=T^7$ removes the inverse letter with constant overhead. Exact arbitrary-angle equivalence is false: the stated ideal rotation gives probability $1/3$, unavailable to the bounded finite-alphabet construction with basis initialization, deterministic classical data and unconditional outputs. Postselection and conditional renormalization are excluded from that obstruction.
\item (primitive 8, optional part) Optional rotations require an effective supply of matrix-entry approximations and a separate approximation theorem. Error is measured on actual matrices, including control-dependent phases; discarding a scalar phase before adding a control is invalid. The general theorem assumes certificates. They are supplied here only for the stated determinant-one one-qubit case and bounded controls, using an inverse-closed determinant-one set whose generated subgroup is dense. The word-length bound is polylogarithmic in inverse error, but the proved certificate search time is exponential in that length bound and includes entry-generation and full rational-request costs.
\end{enumerate}
